\documentclass[12pt,a4paper]{article}
\usepackage[utf8]{inputenc}
\usepackage{amsmath, bm, amssymb, amsfonts}
\usepackage{geometry}
\geometry{margin=2.5cm}
\usepackage{hyperref}
\usepackage{pgfplots}
\usepackage{fancyhdr}
\usepackage{titlesec}

\titleformat{\section}
  {\normalfont\Large\bfseries\centering}  % Format: font size, bold, centered
  {\thesection}{1em}{} 

\pagestyle{fancy}
\pgfplotsset{compat=1.18}
\fancyhead{}
\fancyfoot{}

\title{Analisi Matematica 1 - Esercizi}
\author{}
\date{}

\fancyhead[L]{Gabriele Perna}          % Left side
\fancyhead[C]{Analisi 1 - Esercizi}    % Center
\fancyhead[R]{2025/2026}               % Right side
\fancyfoot[C]{\thepage}                % Page number in the center of footer

\let\oldsection\section
\renewcommand{\section}{\clearpage\oldsection}

\begin{document}

\maketitle
\tableofcontents

\section{Esercitazione 02-10-2025}

\subsection*{Esercizio 1}
La funzione $f : [-1,3] \rightarrow \mathbb{R}$ definita da $f(x) = \sin x + \cos^4 x$
\begin{enumerate}
  \item ha massimo ma non ha minimo
  \item ha minimo ma non ha massimo
  \item \textbf{ha sia massimo che minimo}
  \item è limitata ma non ha né massimo né minimo\\
\end{enumerate}
Il teorema di Weierstrass ci assicura che questa funzione abbia massimo e minimo, in quanto è definita su un intervallo chiuso e limitato.

\subsection*{Esercizio 2}
Sia $A = \{x \in \mathbb{R} : \sin^2 x - 2\sin x < 0\}$. Allora:
\begin{enumerate}
  \item $\inf(A) = 0$
  \item $\bm{ \inf(A) = -\infty}$
  \item $\sup(A) = 2$
  \item $\sup(A) = \sin 2$\\
\end{enumerate}
Partiamo subito col fare la disequazione:
\begin{align*}
  \sin^2 x - 2\sin x &< 0\\
  \sin x (\sin x - 2) &< 0
\end{align*}
La parentesi è sempre negativa, quindi:
\begin{align*}
  \sin x &> 0
\end{align*}
In questo caso conviene aiutarsi con un grafico.\\
Si ha quindi:
\[
A = {x \in \mathbb{R} : \forall k \in \mathbb{Z} \rightarrow 2k\pi < x < 2k\pi + \pi}
\]
Si intuisce che la funzione si estende infinitamente a destra e sinistra.

\subsection*{Esercizio 3}
Sia $f : \mathbb{R} \to \mathbb{R}$ definita da
\[
f(x) = 
\begin{cases}
\dfrac{e^x}{x^2}, & x \neq 0, \\
0, & x = 0.
\end{cases}
\]
Allora:
\begin{enumerate}
  \item $f$ è limitata in $\mathbb{R}$
  \item $\bm{\min\{f(x) : x \in \mathbb{R}\} = 0}$
  \item $f$ è crescente in $\mathbb{R}$
  \item $f$ è dispari\\
\end{enumerate}
Quando $x \ne 0$ la funzione sarà sempre $>0$ in quanto $e^x$ non potrà mai risultare $0$.

\subsection*{Esercizio 4}
L’insieme $\{x \in \mathbb{R} : x^2 > |3x - 1|\}$:
\begin{enumerate}
  \item non è limitato né superiormente né inferiormente
  \item è limitato superiormente ma non inferiormente
  \item è limitato
  \item è limitato inferiormente ma non superiormente
\end{enumerate}

\subsection*{Esercizio 5}
L’insieme $\{x \in \mathbb{R} : x^2 - \tfrac{1}{x} < 0\}$:
\begin{enumerate}
  \item è limitato inferiormente ma non superiormente
  \item è limitato
  \item non è limitato né inferiormente né superiormente
  \item è limitato superiormente ma non inferiormente
\end{enumerate}

\subsection*{Esercizio 6}
L’insieme $\{x \in \mathbb{R} : \tfrac{1}{3} < \sin x \le \tfrac{1}{2}\}$:
\begin{enumerate}
  \item è limitato
  \item è limitato superiormente ma non inferiormente
  \item non è limitato né inferiormente né superiormente
  \item è un intervallo
\end{enumerate}

\subsection*{Esercizio 7}
Sia $A$ l’insieme di definizione della funzione $f(x) = \log(\log(x + 3))$. L’insieme $A$:
\begin{enumerate}
  \item è limitato
  \item non è limitato né superiormente né inferiormente
  \item è limitato superiormente ma non inferiormente
  \item è limitato inferiormente ma non superiormente
\end{enumerate}

\subsection*{Esercizio 8}
La funzione $f : (-2,2) \to \mathbb{R}$ definita da
\[
f(x) = 
\begin{cases}
x + 1, & -2 < x \le 0,\\
x - 1, & 0 < x < 2
\end{cases}
\]
\begin{enumerate}
  \item ha minimo ma non ha massimo
  \item non ha né massimo né minimo
  \item ha sia massimo che minimo
  \item ha massimo ma non ha minimo
\end{enumerate}

\subsection*{Esercizio 9}
L’insieme di definizione della funzione $f(x) = \log(x^2 - 5x + 4 - x)$ è:
\begin{enumerate}
  \item $(-\infty, 4)$
  \item $(-\infty, 0)$
  \item $(-\infty, 0) \cup (4,5)$
  \item $(4,5)$
\end{enumerate}

\subsection*{Esercizio 10}
L’insieme $A = \{x \in \mathbb{R} : |x^2 - 2| > x - 1\}$:
\begin{enumerate}
  \item è limitato superiormente ma non inferiormente
  \item è limitato inferiormente ma non superiormente
  \item non è limitato né superiormente né inferiormente
  \item è limitato
\end{enumerate}

\subsection*{Esercizio 11}
La funzione $f : \mathbb{R} \to \mathbb{R}$ definita da
\[
f(x) = \arctan x + e^{\arctan x}
\]
\begin{enumerate}
  \item è surgettiva ma non iniettiva
  \item è bigettiva
  \item non è né iniettiva né surgettiva
  \item è iniettiva ma non surgettiva
\end{enumerate}
\end{document}
